

\begin{theorem}\label{teo:gauss}
\textbf{Teorema de la divergencia o de Gauss.} Sean $\Omega \subset \mathbb{R}^3$ una región sólida acotada y la superficie cerrada $\partial \Omega$ su frontera,  orientada mediante el vector normal exterior. Sea $A \subseteq \mathbb{R}^3$ un conjunto abierto tal que $\Omega \cup \partial \Omega  \subset A$. Si $\mathbf{F}: A \to \mathbb{R}^3$ es un campo vectorial de clase $\mathcal{C}^1(A)$, entonces:
\[
    \oiint_{\partial \Omega} \mathbf{F} \cdot d\mathbf{S} = \iiint_{\Omega} \nabla \cdot \mathbf{F} \, dV
\]
\end{theorem}



\begin{example}
Calcular el flujo exterior de $\mathbf{F}(x, y, z) = (x, y, z)$ a  trav\'es de $S = \partial \Omega$ siendo   
\[
\Omega = \{(x, y, z) \in \mathbb{R}^3 : x^2 + y^2 + z^2 \leq 1\}.
\]
\end{example}

\noindent\textbf{Solución: } Por el Teorema 10.3.1 (Teorema de la Divergencia de Gauss), tenemos
\[
\iint_S \mathbf{F} \cdot d\mathbf{S} = \iiint_{\Omega} \nabla \cdot \mathbf{F} \, dV.
\]

Sabemos que $F$ es un campo de clase  $\mathcal{C}^1(\mathbb{R}^{3})$  y que $S= \partial \Omega$ es una superficie cerrada. Calculemos la divergencia del campo $\mathbf{F}$:
\[
\nabla \cdot \mathbf{F} = \frac{\partial}{\partial x}(x) + \frac{\partial}{\partial y}(y) + \frac{\partial}{\partial z}(z) = 1 + 1 + 1 = 3.
\] Por lo tanto, aplicando el Teorema de Gauss (\ref{teo:gauss}) queda
\[
\iint_S \mathbf{F} \cdot d\mathbf{S} = \iiint_{\Omega} \nabla \cdot \mathbf{F} \, dV = \iiint_{\Omega} 3 \, dV = 3 \cdot \text{Vol}(\Omega) = 3 \cdot \frac{4}{3}\pi = 4\pi.
\]

En consecuencia, el flujo saliente de $\mathbf{F}$ a trav\'es de $S$ es
\[
\boxed{\iint_S \mathbf{F} \cdot d\mathbf{S} = 4\pi}.
\]



\begin{tarea}
 Probar que el flujo de $\mathbf{F} = (ax, by, cz)$ con $a,b,c \in \mathbf{R}^+$ a través del cubo: $V : [-1,1] \times [-1,1] \times [-1,1]$ orientado con normales salientes es $(a+b+c) \mathrm{volumen}(V)$.
\end{tarea} 


\begin{tarea}
 Sea la superficie   $S=\{  (x,y,z) \in \mathbb{R}^{3} \: : \: x ^{2}+ y^{2} = z^{2}, 1 \leq z \leq 2 \}$   y   sea el campo $F(x,y,z)=(x+y^{2}z^{2}, y+z^{2}x^{2}, z).$  Hallar el flujo  de $F$ a trav\'es de $S$ indicando la orientaci\'on elegida.
\end{tarea}  



\begin{tarea}
Sea  $S=\{  (x,y,z) \in \mathbb{R}^{3} \: : \: x ^{2}+ y^{2} +1= z,  \:  1\leq z \leq h\}$  y sea  el campo $F(x,y,z)=(z+y^{3},  y+ \cos (z^{2}x^{2}),  -z).$    Hallar $h \in \mathbb{R}$ tal que: $$\int_{S} F . dS = 2\pi. $$  Indicar la orientaci\'on con la cual resuelve el ejercicio.
\end{tarea}